% Generated by min_heterogeneity.R
\begin{table}[htbp]
\centering
\footnotesize
\setlength{\tabcolsep}{3pt}
\caption{Tax rates minimising the heterogeneity of costs across Member States.}
\label{tab:min_heterogeneity}
\begin{tabular}{@{}l*{9}{r}@{}}
\toprule
 & Report & \multicolumn{4}{c}{Uniform luxury VAT} & \multicolumn{4}{c}{Luxury taxed by category} \\
\cmidrule(lr){3-6}\cmidrule(lr){7-10}
 & rates & \makecell{Sq.\\dev.} & Minimax & Frugal & \makecell{Max.\\gain} & \makecell{Sq.\\dev.} & \textbf{Minimax} & Frugal & \makecell{Max.\\gain} \\
\midrule
\multicolumn{10}{@{}l}{\textit{Tax rates (\%)}} \\
DST & 3 & 12 & 7.7 & 12 & 12 & 12 & \textbf{6.3} & 12 & 12 \\
Ad tax & 15 & 50 & 60 & 0 & 60 & 36 & \textbf{60} & 18 & 60 \\
FTT & .5 & .59 & .63 & .49 & 1.7 & .74 & \textbf{.56} & .38 & 1.7 \\
Wealth tax & .5 & .51 & .66 & .99 & 0 & 0 & \textbf{.49} & .27 & 0 \\
Aviation tax & 20 & 4.9 & 6.9 & 36 & 24 & 2.7 & \textbf{3.1} & 0 & 24 \\
Luxury tax & 20 & 3.7 & 0 & 0 & 0 &  &  &  &  \\
\quad automotive & 20 &  &  &  &  & 40 & \textbf{5.4} & 35 & 0 \\
\quad yachts & 20 &  &  &  &  & 15 & \textbf{33} & 29 & 0 \\
\quad design \& furniture & 20 &  &  &  &  & 80 & \textbf{80} & 80 & 0 \\
\quad gourmet \& dining & 20 &  &  &  &  & 65 & \textbf{0} & 80 & 0 \\
\quad personal goods & 20 &  &  &  &  & 0 & \textbf{0} & 5.9 & 0 \\
\quad hospitality & 20 &  &  &  &  & 0 & \textbf{0} & 0 & 0 \\
\quad wines \& spirits & 20 &  &  &  &  & 0 & \textbf{0} & 0 & 0 \\
\midrule
RMSD from mean (\% GNI) & .31 & .152 & .183 & .219 & .212 & .12 & \textbf{.169} & .141 & .212 \\
Max.\ cost (\% GNI) & 1.628 & 1.02 & .884 & 1.266 & 1.323 & .912 & \textbf{.852} & 1.039 & 1.323 \\
Highest-cost country & CY & CY & LU & CY & NL & CZ & \textbf{LU} & LU & NL \\
Highest-cost frugal country & AT (.75) & NL (.85) & NL (.87) & AT (.7) & NL (1.32) & NL (.9) & \textbf{NL (.85)} & AT (.7) & NL (1.32) \\
New funds (EUR bn) & 29.8 & 18.8 & 20.9 & 31.1 & 0 & 14 & \textbf{20.5} & 24.1 & 0 \\
\bottomrule
\end{tabular}
\begin{minipage}{\textwidth}\vspace{4pt}\footnotesize \textit{Note:} Total new revenues are held at \EUR 103 bn (.705\% of EU GNI), and each rate is capped at 4 times the report's rate. The criteria are: minimising the sum over the Member States of squared deviations of costs from the EU average (Sq.\ dev.), minimising the maximum cost (Minimax), minimising Sq.\ dev. under the constraint that the frugal countries (AT, DE, IE, NL, SE) pay less than the EU average (Frugal), and maximising the minimum State budget gain (Max.\ gain).\end{minipage}
\end{table}
